\scanpage{48}

\uline{Questions sur} V.A. \uline{et groupes formels}.

1) (Barsotti). Soient $A,B$ des VA sur $k$ alg clos (car $p>0$) ``hypersingulières'' i.e. sans pts d'ordre $p$. Supposons que les groupes formels $\overline A,\overline B$ soient \uline{isogènes}. Alors $A$ et $B$ sont-elles isogènes ?

\inserted{Si \deleted{\unsure{p}} $\dim A=1$, c'est connu dans Deuring.}

Manin ignore déjà la réponse si $A,B$ de dim $2$ (donc $\overline A,\overline B$ sont \uline{isogènes} : $G_{1,1}+G_{1,1}$). La question revient alors à ceci Si $A$ de dim $2$ est hypersingulière, est elle isogène à un produit de deux courbes elliptiques ($\Leftrightarrow$ à un carré de courbe elliptique) --- nécessairement hypersingulière. Il affirme que c'est vrai en car $p=2$.

Noter que si on abandonne l'hypothèse supersingulière, et qu'on suppose $A,B$ \deleted{\unsure{simples}} de dim $2$, à groupes formels isomorphes non toroïdaux il n'est pas vrai en général que $A$ et $B$ isogènes \inserted{$\overline{\mathrm m}$ si $A$ et $B$ simples.} En effet, d'après Manin, les \inserted{\uline{jacobiennes}} $A$ de dim $2$ \deleted{(polarisées principales)} de grp formel $\widehat G_m+G_{1,1}$ forment une sous famille \inserted{$N$} \unsure{loc. fermée} de \uline{dim $2$} de la famille modulaire \inserted{$M_2$}. Parmi celles-ci, celles qui sont \deleted{isogènes} \inserted{non} simples i.e. \deleted{\unsure{décomposées}} isogènes à un produit $E\times F$, $E$ et $F$ courbes elliptiques dont l'une \inserted{(soit $F$)} hypersingulière l'autre non, forment une famille de dim $1$ au plus. \margin{[ \uline{NB.} il n'y a pas de famille \unsure{continue}}

\margin{58}
